\section{1re loi modulaire pour les variations de l’aire}

Définissons les deux aires normalisées
\begin{equation}
\label{rm:eq:modular-normalized-area-channels}
\mathcal A_m=\gammaH a_0N_m,
\qquad m\in\{90,120\}.
\end{equation}
Alors
\begin{equation}
\label{rm:eq:modular-beta-definition}
\HHM^{(A)}-cI
=\beta_{90}\mathcal A_{90}+\beta_{120}\mathcal A_{120},
\qquad
\beta_m=\frac{m\Dmod}{\gammaH a_0}.
\end{equation}
Les valeurs terminales sont
\begin{align}
\beta_{90}&=5.9726485401070929914651798889\ldots,
\label{rm:eq:modular-beta90}\\
\beta_{120}&=7.9635313868094573219535731852\ldots,
\qquad
\frac{\beta_{120}}{\beta_{90}}=\frac43.
\label{rm:eq:modular-beta120}
\end{align}

\begin{theorem}[Première loi modulaire en variables aréales]
\label{rm:thm:modular-first-law}
Pour une famille différentiable d’états normalisés \(\rho(\varepsilon)\) passant par \(\rho_A\), tout en maintenant fixes \(\Dmod,\gammaH,a_0\),
\begin{equation}
\label{rm:eq:modular-first-law}
\boxed{
\delta S
=\beta_{90}\,\delta\langle\mathcal A_{90}\rangle
+\beta_{120}\,\delta\langle\mathcal A_{120}\rangle.}
\end{equation}
\end{theorem}

\begin{proof}
La première loi d’intrication dans l’état de référence est \(\delta S=\delta\langle-\log\rho_A\rangle\). Le terme scalaire \(cI\) ne contribue pas puisque \(\delta\Tr\rho=0\). Substituer \cref{rm:eq:modular-beta-definition} prouve l’identité.
\end{proof}

La formule est locale dans l’espace des états. Pour des changements finis, l’entropie relative apparaît ; l’omettre transformerait une identité tangente en une assertion fausse de linéarité globale.

\Needspace{22\baselineskip}
\begin{theorem}[Identité finie avec terme de défaut]
\label{rm:thm:modular-finite-first-law}
Pour tout état \(\sigma\) dont le support est contenu dans celui de \(\rho_A\),
\begin{equation}
\label{rm:eq:modular-finite-first-law}
\boxed{
S(\sigma)-S(\rho_A)
=\beta_{90}\Delta_\sigma\langle\mathcal A_{90}\rangle
+\beta_{120}\Delta_\sigma\langle\mathcal A_{120}\rangle
-D(\sigma\Vert\rho_A),}
\end{equation}
où \(D(\sigma\Vert\rho_A)=\Tr\sigma(\log\sigma-\log\rho_A)\geq0\).
En particulier, la variation linéaire de \cref{rm:thm:modular-first-law} est le terme tangent de cette identité.
\end{theorem}

\begin{proof}
Par définition et par \(\HHM^{(A)}=-\log\rho_A\),
\[
D(\sigma\Vert\rho_A)
=-S(\sigma)+\Tr(\sigma\HHM^{(A)}).
\]
Pour \(\sigma=\rho_A\), le membre de gauche est nul et \(S(\rho_A)=\Tr(\rho_A\HHM^{(A)})\). La soustraction des deux identités donne
\[
S(\sigma)-S(\rho_A)
=\Delta_\sigma\langle\HHM^{(A)}\rangle
-D(\sigma\Vert\rho_A).
\]
Le scalaire \(cI\) s’annule entre états normalisés ; substituer \cref{rm:eq:modular-beta-definition} produit \cref{rm:eq:modular-finite-first-law}. La positivité de l’entropie relative est l’inégalité de Klein.
\end{proof}

La correction \(D(\sigma\Vert\rho_A)\) mesure la distance informationnelle entre l’état réel et l’état modulaire de référence. Ainsi, la première loi ne perd pas son contenu au-delà du régime infinitésimal : elle devient une identité exacte avec un défaut positif contrôlé.
